\section{Mathematische Definitionen und Begriffe}

Im Folgenden verwenden wir den Begriff \textit{Abbildung} synonym zu dem Begriff \textit{Funktion}.

\section*{A}

\begin{description}
    \item[Abbildung]
    Siehe \textit{Funktion}.
\end{description}

\section*{B}

\begin{description}
    \item[Backus-Naur-Form]
\end{description}
\begin{description}
    \item[Bottom Typ \bot]
    Der Bottom-Typ ist bezeichnet in der Typtheorie einen Typen, vom dem aus keine weiteren Typen mehr abgeleitet werden können - oder im Zusammenhang mit \textit{Top Typ} -- der als Subtyp aller anderen Typen existiert.
    Der Typ wird auch als \textit{leerer Typ} bezeichnet und wird in der Regel zur formalen Darstellung von \texttt{nullptr} in C++ oder \texttt{NULL} in Java verwendet.
\end{description}

\section*{D}

\begin{description}
    \item[Disjunkte Vereinigung]
    Die \textit{disjunkte Vereinigung}
    \[
        A \dot\cup B
    \]
    \noindent
    bezeichnet die Vereinigung zweier paarweise disjunkter Mengen.
    Die Notation setzt also voraus:
    \[
        A \cap B = \emptyset
    \]
\end{description}

\section*{E}
\begin{description}
    \item[Exklusion]
    Die \textit{Exklusion}
    \[
        A \backslash B
    \]
    \noindent
    ist die Menge, für die gilt
    \[
        a \in A \land a \notin B
    \]
\end{description}

\section*{F}

\begin{description}
    \item[Funktion] Eine Funktion ordnet einem Element eines Definitionsbereichs $D$ genau ein Element des Wertebereichs $W$ zu.
    \[
        f: D \rightarrow W
    \]
    Eine Funktion ist, sofern nichts anderes vereinbart, immer \textit{total}.
\end{description}


\begin{description}
    \item[Funktionale Abhängigkeit]
\end{description}

\section*{H}

\begin{description}
    \item[Hülle (Funktionale Abhängigkeit)]
\end{description}


\section*{J}

\begin{description}
    \item[JOIN Abhängigkeit]
\end{description}


\section*{K}

\begin{description}
    \item[Kontextfreie Grammatik]
    Eine Grammatik $G$, bestehend aus einer Menge von Terminalsymbolen $\Sigma$, Nichtterminalsymbolen $N$, Produktionen bzw. Regeln $P$ und einem Startsymbol $S$
    \[
        G = (\Sigma, N, P, S)
    \]
    \noindent
    heißt kontextfrei, wenn falls ihre Regelmenge
    \[
        P \subseteq N \times (\Sigma \cup N)^*
    \]
    \noindent
    ist und außerdem
    \[
      \mid P \mid < \infty
    \]
    \noindent
    erfüllt.
    Der Begriff \textit{kontextfrei} beschreibt die Eigenschaft, dass die linke Seite einer Produktion bei der Ableitung unabhängig von dem Kontext des bereits abgeleiteten Wortes ist.\\
    \noindent
    Bei \textit{kontextsensitiven Grammatiken} ist die linke Seite einer Produktion abhängig von einem Kontext, das zu ersetzende Nichtterminal taucht also in Verbindung mit anderen (Nicht)Terminalsymbolen auf.
 \end{description}


\section*{K}

\begin{description}
    \item[Minimale Funktionale Abhängigkeitsmenge]
\end{description}

\section*{P}

\begin{description}
    \item[Partielle Funktion]
    Für \textit{partielle Funktionen} notieren wir $f : D \rightharpoonup W$.
    Eine Funktion $f$ ist \textit{partiell}, wenn sie für Elemente $x \in D$ \textit{undefiniert} sein kann.
    Gibt es ein solches $x$, dann notieren wir in diesem Fall
    \[
        f(x) = \bot
    \]
    \noindent
    Deshalb sind partielle Funktionen nicht links-tota, aber die Eigenschaft der \textit{Rechtseindeutigkeit} verbleibt, da jedem $x \in D$ höchstens ein $f(x) \in W$ zugeordnet ist.
\end{description}

\section*{R}

\begin{description}
    \item[Rechtseindeutige Relation]
    Eine Relation $R \subseteq A \times B$ ist \textit{rechtseindeutig}, wenn jedem Element $x \in A$ höchstens ein Element $y \in B$ zugeordnet ist:
    \[
        \forall x \in A, y_1, y_2 \in B : (x, y_1) \in R \land (x, y_2) \in R \Rightarrow  y_1 = y_2
    \]
\end{description}


\begin{description}
    \item[Relation (Datenbank)]
\end{description}

\begin{description}
    \item[Relationsschema (Datenbank)]
\end{description}


\begin{description}
    \item[Relationszustand (Datenbank)]
\end{description}


\section*{S}

\begin{description}
    \item[Superschlüssel]
\end{description}

\section*{T}

\begin{description}
    \item[Top Typ]
    Der \textit{Top Typ} ist in der Typ-Theorie der Supertyp aller in der darunterliegenden Programmiersprache vorkommenden Typen, beispielsweise \texttt{Object} in Java (siehe auch \textit{Bottom Typ}).
\end{description}

\begin{description}
    \item[Totale Funktion]
    Für \textit{totale Funktionen} definieren wir $f: D \rightarrow W$.
    Eine Funktion $f$ ist \textit{total}, wenn gilt:
    \[
        \forall x \in D: f(x) \in W
    \]
    \noindent
    Eine totale Funktion ist links-total und rechtseindeutig.
\end{description}

\section*{U}

\begin{description}
    \item[Undefiniert]
    Für einen \textit{undefinierten Wert} vereinbaren wir das Symbol $\bot$.
\end{description}

\begin{description}
    \item[Undefinierter Zustand]
    Ein undefinierter Zustand ist im Kontext dieser Arbeit ein Zustand, dessen Auftreten zu einem Abbruch der weiteren VErarbeitung führt, bzw, in einem Modell zu einer allgemeine ungültigen Bedingung fürht, für die das Modell nciht definiert ist.
\end{description}

\begin{description}
    \item[Unit Type \top]
    Der \textit{Unit Type} $\top$ ist ein Basistyp in Metasprachen wie ML.
    In modernen Programmiersprachen ist er zu vergleichen mit dem Typ \texttt{void}, beispielsweise als
    Rückgabetyp von Funktionen.\footnote{
        Pierce weist in diesem Kontext darauf hin, dass dem Typ eine semantisch andere Rolle zukommt als dem \textit{leeren} Typen \bot.
    }
\end{description}

\section*{V}


\begin{description}
    \item[Verlustfreier Join]
\end{description}